\scanpage{40}

% B1 | text | opening paragraph
Dans le cas qui nous intéresse ($\unsure{\xi|\zeta}\in\mathrm{Br}(\mathbb Q)$
d'indice $\unsure{n=2}$, ramifié \unsure{seulmt} en $2$ et $\infty$). \deleted{on trouve}
on trouve que $\unsure{\mathrm{Bord}|\mathrm{Brd}}(C)$ est la \deleted{\unsure{gr.}}\inserted{catégorie} de Picard
\unsure{splittée} d'invariants $\mathbb Z/2\mathbb Z$, $\mathbb Z/2\mathbb Z$, le générateur
$\unsure{\beta|\rho}$ de $\pi_0$ laquelle correspond au \deleted{\unsure{l'endomorphisme}} foncteur
qui à une courbe elliptique $E$ associe $E^{(p)}=E/\operatorname{Ker} F_E$)
correspond à l'unique idéal premier \inserted{$\mathfrak m$} de $O$ sur $p$,

% B2 | display | norm and choice in upper third
\[
\text{et on a}\qquad \underset{\substack{\Vert\\p\mathbb Z}}{N_{O/\mathbb Z}(\mathfrak m)}
\hookrightarrow N_{O/\mathbb Z}(O)=\mathbb Z,\qquad \text{on choisit }\bar\beta=p.
\]

% B3 | text | middle, first conditional
Ceci nous dit que \uline{si} $p$ \uline{est inversible dans} $A'$, alors
l'image de $(\beta,\bar\beta)$ dans $\mathrm{Brd}(C_{A'},T_{A'})\simeq\deleted{\unsure{\mathrm{Brd}(C_{A'})}}H^1(S',\mu_2)$
est la classe de $p$ modulo des carrés.

% B4 | text | middle, local case, with marginal notation
Par contre, prenons le cas où $A'=\mathbb Z_p$, alors
\deleted{on est dans un groupe}
$\mathrm{Brd}(C_p,T_p)\simeq\mathrm{Brd}(C_{A'},T_{A'})$ est une extension de $\mathbb Z/2\mathbb Z$ par
\margin{$C_p=C'$\newline $T_p=T'$}
$\mathbb Z_p^*/\mathbb Z_p^{*2}$ ($\simeq\mathbb Z/2\mathbb Z$ si $p\ne2$), en fait c'est
$\underbrace{\mathbb Z_p^*/\mathbb Z_p^{*2}}_{=H^1(\mathbb Z_p,\mu_2)\deleted{\unsure{\mathbb Z/(2)}}}\times
\underbrace{\mathbb Z/2\mathbb Z}_{\substack{\text{engendré par}\\(\mathfrak m,p)}}$, et les
classes d'isom \inserted{possibles des} données à l'$\infty$ \deleted{\unsure{liées aux}} forment un torseur
sous ce groupe (\deleted{\unsure{le plus souvent}} $\simeq\mathbb Z/2\mathbb Z\times\mathbb Z/2\mathbb Z$ si $p\ne2$).

% B5 | text | last paragraph, down to foot
D'autre part, si $A'=\mathbb R$, on peut considérer les classes d'isomorphie de la gerbe de lien $\mu_n$ sur $\mathbb R$ définie par $C$
(soit $\unsure{\mathscr G|\mathcal G}_{\mathbb R}$, où $\mathscr G$ est la gerbe sur $\mathbb Z[\frac1p]$) et \deleted{l'ensemble} des
classes d'\inserted{isomorphie d'}équivalences avec la gerbe-type sur $\mathbb R$
($=$ ensemble des \unsure{\uline{ple}ssions|\uline{pl}sections} de la gerbe), qui est un
torseur sous $H^1(\mathbb R,\mu_2)=\mathbb R^*/\mathbb R^{*2}$. On trouve donc
% archive sheet number: 39
